% Generated draft: requires mathematical and source review.
\documentclass[11pt]{article}
\usepackage[margin=25mm]{geometry}
\usepackage{fontspec}
\setmainfont{Latin Modern Roman}
\usepackage{amsmath,amssymb,mathtools,unicode-math}
\setmathfont{Latin Modern Math}
\usepackage{enumitem,needspace,xurl,hyperref}
\setlength{\parindent}{0pt}
\setlength{\parskip}{5pt}
\setlength{\emergencystretch}{3em}
\newcommand{\N}{\mathbb{N}}
\newcommand{\Z}{\mathbb{Z}}
\newcommand{\Q}{\mathbb{Q}}
\newcommand{\R}{\mathbb{R}}
\newcommand{\C}{\mathbb{C}}
\newcommand{\F}{\mathbb{F}}
\newcommand{\A}{\mathbb{A}}
\newcommand{\PP}{\mathbb{P}}
\newcommand{\E}{\mathbb{E}}
\newcommand{\eps}{\varepsilon}
\newcommand{\dd}{\,\mathrm{d}}
\begin{document}
\section*{UnsolvedMath ID: 30005513}
\textbf{Source problem code:} catalogue entry
\subsection{Definitions and mathematical statement}
Let \(\ell\) be a prime, let \(G\) be a finite group, and, under the usual modular representation theory convention, let \(k\) be an algebraically closed field of characteristic \(\ell\). An \(\ell\)-block \(B\) of \(kG\) is a direct summand \(kGb\), where \(b\) is a primitive central idempotent of \(kG\). The supplied statement does not specify the hypotheses on \(k\); the convention just stated makes the intended block fusion system precise and requires confirmation against the source.

Choose a maximal \(B\)-Brauer pair \((D,e_D)\); its underlying \(\ell\)-subgroup \(D\) is a defect group of \(B\). Here a \(B\)-Brauer pair is a pair \((Q,e_Q)\), where \(Q\leq G\) is an \(\ell\)-subgroup, \(e_Q\) is a block idempotent of \(kC_G(Q)\), and \(\operatorname{Br}_Q(b)e_Q\neq 0\), with \(\operatorname{Br}_Q\) the Brauer homomorphism. For each \(Q\leq D\), let \((Q,e_Q)\) be the unique \(B\)-Brauer pair contained in \((D,e_D)\), using the usual partial order on Brauer pairs.

The block fusion system \(\mathcal F_D(B)\) has all subgroups of \(D\) as objects and morphism sets
\[
\operatorname{Hom}_{\mathcal F_D(B)}(Q,R)
=\bigl\{c_g|_Q:Q\longrightarrow R\;\bigm|\;
 g\in G,\ {}^g(Q,e_Q)\leq(R,e_R)\bigr\},
\qquad Q,R\leq D,
\]
where \(c_g(x)=gxg^{-1}\) and \({}^g(Q,e_Q)=(gQg^{-1},ge_Qg^{-1})\). The supplied source states that this fusion system is saturated.

For a finite group \(H\) and a Sylow \(\ell\)-subgroup \(S\leq H\), the group fusion system \(\mathcal F_S(H)\) has subgroups of \(S\) as objects and
\[
\operatorname{Hom}_{\mathcal F_S(H)}(U,V)
=\{c_h|_U:U\longrightarrow V\mid h\in H,\ hUh^{-1}\leq V\}.
\]
A saturated fusion system is called \emph{nonexotic} if it is isomorphic to \(\mathcal F_S(H)\) for some such finite group \(H\) and Sylow subgroup \(S\).

\textbf{Statement to be proved.} For every prime \(\ell\), every finite group \(G\), and every \(\ell\)-block \(B\) as above, the saturated fusion system \(\mathcal F_D(B)\) is nonexotic. Explicitly, there exist a finite group \(H\), a Sylow \(\ell\)-subgroup \(S\) of \(H\), and an isomorphism \(\alpha:D\longrightarrow S\) such that, for every \(Q,R\leq D\),
\[
\operatorname{Hom}_{\mathcal F_S(H)}(\alpha(Q),\alpha(R))
=\bigl\{\alpha\circ\varphi\circ(\alpha|_Q)^{-1}
\mid \varphi\in\operatorname{Hom}_{\mathcal F_D(B)}(Q,R)\bigr\}.
\]
The realizing group \(H\) may depend on \(B\). The supplied registry labels this problem open; that label is supplied curation, not an independently verified mathematical or literature assessment.
\subsection{Short English statement}
Every block of a finite group determines a saturated fusion system on a defect group. The problem asks whether this fusion can always be realized by conjugation in some finite group in which the defect group is a Sylow subgroup. The coefficient field convention needs source confirmation, and the supplied open status requires independent review.
\subsection{Sources}
\begin{itemize}
\item Local registry: \texttt{data/unsolvedmath/problems.json}, permanent UnsolvedMath ID \(30005513\), title ``Nonexoticity of Block Fusion Systems,'' exact source code \texttt{catalogue entry}. The supplied record gives proposed year 2023 and no named proposer.
\item Supplied primary report citation: \emph{Representations of Finite Groups} (2023), Oberwolfach Reports, DOI 10.4171/owr/2023/19, pp. 1031--1084. \url{https://doi.org/10.4171/owr/2023/19}. The registry marks the statement \texttt{corrected\_verified} and reports reconstruction and checking against the official report; this draft does not independently verify that assertion.
\item Supplied literature source: Serwene, ``Block-exoticity of a family of exotic fusion systems.'' \url{https://doi.org/10.1515/jgth-2019-0035}.
\item Supplied literature source: Serwene, ``proving a conjecture for fusion systems.'' \url{https://arxiv.org/abs/2302.05203}.
\item Supplied literature source: Oberwolfach Report 19/2023. \url{https://doi.org/10.4171/owr/2023/19}.
\item Literature assessment dated exactly \texttt{2026-08-22}: the supplied curation labels the problem open, states that block-exoticity has been proved for several families, and states that no general realization theorem is known. This is supplied dated curation derived from the OWR record, not independently verified here; time-sensitive status and cited claims require independent checking.
\item The raw excerpt also contains the heading ``Spetses from the p-local perspective,'' naming Jason Semeraro and joint work with Radha Kessar and Gunter Malle. Its relationship to the problem's attribution is not established by the supplied extraction.
\end{itemize}
\end{document}
